\section{Introduction}
\label{sec:intro}

3D Gaussian Splatting (3DGS) represents a scene as a set of 3D Gaussians and renders it fast enough for interactive viewing~\cite{kerbl2023gaussian}.
The original method fits the Gaussians to each scene by optimization, which takes minutes and needs the whole capture.
Feed-forward models remove this step: a trained network predicts the Gaussians from the input frames in one pass~\cite{anysplat2025,streamsplat2026}.
StreamSplat reads a monocular video one keyframe at a time and outputs Gaussians with their motion~\cite{streamsplat2026}.
On our GPU, one StreamSplat update costs about 0.19~s with the released code. % log: G5
At this cost, a server can keep a live 3D model of a user's camera stream and refresh it about once per second. % log: G5, G6

Many users on one GPU make this a GPU virtualization problem.
Each user is a tenant with its own model state and deadline.
The server must know how many tenants it can admit, how fast each update returns, and how much quality it can trade for speed.
To our knowledge, no published system serves many users of a feed-forward 3DGS model from one GPU (Section~\ref{sec:related}).

We propose a framework for this setting with three goals: serve more users, lower each user's latency, and keep the reconstruction quality.
The framework compares two levers (Figure~\ref{fig:system}).
The first lever is the sharing mechanism: time-slicing, NVIDIA's Multi-Process Service (MPS), or batching across users.
These mechanisms change where the work of several users runs.
The second lever is the work per update: a feature cache, bfloat16 (bf16) arithmetic, CUDA graphs and a bf16 depth model.
These changes shrink the work of every update.
We measure capacity, the largest number of users that meet a latency service-level objective (SLO), on one NVIDIA RTX 6000 Ada.

The sharing mechanisms did not add users to the released code.
MPS served as many users as time-slicing in all four settings we ran. % log: G2, G3, G6, G7
Batching eight users' updates into one call gave 0.92 times the throughput of one update per call. % log: G9
Less work per update did add users.
At one keyframe per user per second, the released code serves 4 users, the optimized server 10, and the server with a bf16 depth model 14. % log: G15, G16
The cost is 0.19~dB of keyframe PSNR, and 14 users use 4,884~MiB, about 10\% of the card's memory. % log: G16

This report makes three contributions.
\begin{itemize}
\item A framework and harness for multi-user feed-forward 3DGS serving.
It adapts SLO-based capacity from DNN serving~\cite{clockwork2020} to stateful per-user streams, and it times each optimization against an interleaved matched control.
\item A negative result for the sharing mechanisms on the released code: time-slicing and MPS add no users, and batching does not raise throughput.
\item An ablation of five per-update optimizations that raises capacity from 4 to 14 users and lowers the worst p95 latency, with its quality cost measured on DAVIS 2017 val. % log: G11-G16
\end{itemize}
This progress report covers the server side.
It measures capacity and latency only, not the allocation of GPU time between unequal users or their isolation.
Section~\ref{sec:continuation} describes the browser side and the next experiments.
